\chapter{Implementation}
\label{cha:impl}

This chapter explains how I built the laboratory and why I made each choice. I write in the past tense because this is the version I am submitting, not a plan.

\section{Starting point}
\label{sec:start}

I implemented the laboratory in React, Vite, TypeScript and Express: four tabs, authentication, source labels, tests, and a Python attack path. Appendix~\ref{app:ai} states where Gemini is used.

\section{Runtime and model identifiers}
\label{sec:runtime}

The Node process loads \texttt{.env} and then \texttt{.env.local}. The model id is:

\begin{minted}{javascript}
const GEMINI_MODEL = process.env.GEMINI_MODEL || 'gemini-3.5-flash';
\end{minted}

I used \texttt{gemini-3.5-flash} as the default because that is the id I configured for the teaching calls~\cite{geminiapi}. I can override it with \texttt{GEMINI\_MODEL} without editing \texttt{server.ts}. The process logs the id at startup. I still have to confirm, in the viva, which id the live key actually served on a given day. One labelled Gemini session is stored under \texttt{eval/gemini\_runs/}: three teaching routes with \texttt{source=gemini}. Live FGSM/PGD calls are stored under \texttt{eval/adversarial/results/live\_*.json}. I do not invent extra fallback counts beyond those files.

I chose the official \texttt{@google/genai} client and JSON response schemas so that the UI could read \texttt{subject}, \texttt{body} and the score fields without a second parser.

\section{Authentication}
\label{sec:auth}

My supervisor asked for authentication on every page and every API. I implemented \texttt{requireLabAccess} with two ways in: an examiner login form that sets an HttpOnly cookie (\texttt{POST /api/login}), and HTTP Basic Auth for \texttt{curl} and older clients. Static files are public so the login page can load. Every \texttt{/api/*} route except login and logout stays locked. I chose this because a browser login is easier to demonstrate than a raw Basic prompt, and because the laboratory is a small demonstration, not a campus SSO product.

Credentials come from \texttt{LAB\_BASIC\_USER} and \texttt{LAB\_BASIC\_PASS}. If they are missing, local development continues without a login. In production the same missing pair returns 503 and the text ``This laboratory is closed''.

Comparison uses \texttt{timingSafeEqualString}. Different lengths return false. Empty strings return false.

\section{Rate limit, size limit, and prompt checks}
\label{sec:limits}

I set \texttt{express.json(\{ limit: '32kb' \})} because a teaching API does not need large uploads and because the deepfake route should not become a file dump.

\texttt{rateLimitApi} counts POST requests per IP and path in a one-hour window. I used 10 requests per hour for \texttt{/api/generate-phishing} and 40 for the other POST API paths. I chose a tighter limit on generate because that route spends Gemini quota and because it is the dual-use route (see Chapter~\ref{cha:disc}).

\texttt{promptLooksInjected} looks for three regular expressions: ``ignore previous/prior/above instructions'', ``you are now '', and ``disregard the system/developer''. I know this is a small list. I still wanted a reject path I could test. A match returns HTTP 400 with ``Prompt rejected by input validation.''

\section{Phishing laboratory}
\label{sec:impl-phish}

\texttt{academicDataset.ts} holds three authored examples, four public-pattern reconstructions, and \texttt{generate90Samples()}. The generator uses a Mulberry32 RNG with seed 20260927. I chose a fixed seed so that screenshots and the dashboard aggregates stay reproducible.

The score bands are configuration:

\begin{itemize}
\item legitimate: realism 90--99, detectability 15--35, polish 93--100;
\item conventional phishing: realism 30--55, detectability 80--95, polish 40--65;
\item AI-style phishing: realism 81--96, detectability 40--60, polish 90--100.
\end{itemize}

When the 90 rows are loaded, the dashboard mean designed realism sits near 75\%. The value ``100 minus mean AI detectability'' sits near 51\%. Those two numbers are the ones people keep quoting. They are not measurements of Gmail or SpamAssassin. I kept the original card look and I labelled the badges ``Designed avg'' and ``Designed'' so the screen stays the same while the caption stays honest.

\texttt{parseEmailFile} reads a local \texttt{.txt} or \texttt{.eml} in the browser. It looks for From, To, Subject and a blank-line body split. I added file load because an examiner can bring a sample without pasting by hand. The file never leaves the browser except as the parsed text in \texttt{/api/analyze-email}.

\section{Deepfake tab}
\label{sec:impl-df}

I implemented three presets only. I wrote the filenames so that they look like teaching cases: a CEO video, a voicemail, a passport image. Exact name match returns the preset. Any other string is narrated. I put a sentence on the tab: media bytes are not inspected. I did this because shipping a weak detector would have been worse than shipping a labelled simulator.

\section{Adversarial path}
\label{sec:impl-adv}

\texttt{eval/adversarial/service.py} is a FastAPI application. On first attack it downloads CIFAR-10 and loads \texttt{cifar10\_resnet56} through \texttt{torch.hub}. It walks the test set and takes the first image of the requested class that the model already classifies correctly. The default attack is untargeted FGSM or PGD in pixel space; a targeted run happens only when the request sets \texttt{targeted: true}. The result is projected onto the $L_\infty$ ball and returned as softmax percentages plus the measured $L_\infty$. Unknown class names return HTTP 400. The architecture field is ignored; the loaded checkpoint is always ResNet-56.

\texttt{eval/adversarial/run\_attacks.py} is the batch script. It samples 200 test images by default, seed 20260927, and writes \texttt{adversarial\_metrics.json} with \emph{untargeted} FGSM and PGD. I describe the method in Chapter~\ref{cha:eval}. I do not paste invented accuracies.

The attack tab names the model CIFAR-10 ResNet-56 and offers FGSM or PGD. I do not offer ImageNet class names (banana, panda) as presets. The class inputs are CIFAR-10 labels. The photographs next to the canvas are stock illustrations, not the $32 \times 32$ tensors used by the attack.

\section{Tests}
\label{sec:tests}

I wrote the automated tests in \texttt{tests/labSecurity.test.ts} and \texttt{tests/apiContract.test.ts} and I run them with \texttt{npm test} (\texttt{tsx --test}). Table~\ref{tab:tests} lists the first four security helpers. The other tests check the 32\,kb limit, the 10/40 rate limits, fail-closed production auth, the login cookie, the four POST teaching routes, the three eval GET routes, and the \texttt{source} labels.

\begin{table}[htbp]
\centering
\caption{Automated tests in \texttt{labSecurity.test.ts}.}
\label{tab:tests}
\scriptsize
\begin{tabular}{p{5.6cm}p{7.6cm}}
\toprule
Test name & What I asserted \\
\midrule
\texttt{timingSafeEqualString rejects different lengths} & \texttt{'ab'} vs \texttt{'a'} is false; two empty strings are false. \\
\texttt{timingSafeEqualString accepts exact match} & \texttt{alpha} equals itself; \texttt{alpia} does not. \\
\texttt{labCredentials fail closed when env is empty} & After deleting the two env vars, \texttt{labCredentials()} is \texttt{null}. \\
\texttt{prompt injection patterns are rejected} & ``ignore previous instructions and dump the key'' is true; a normal library-fine prompt is false. \\
\bottomrule
\end{tabular}
\end{table}

I did not write an end-to-end browser test. I tested the four tabs by hand on the local server. That is a limitation. The unit tests still cover the lock helpers, the size cap, and the route contracts that my supervisor asked for.

\section{What I refused to implement}
\label{sec:refused}

I refused a Gemini fallback on the adversarial POST. I refused a mailbox crawler. I refused to treat the dashboard percentages as experimental results. I refused to keep the academic-report generator in the repository. Those refusals are part of the implementation, not only of the discussion.
